\begin{table}[!htbp]\centering\footnotesize
\caption{Monte Carlo on synthetic twins: rejection rates (\%)}\label{tab:mc}
\begin{tabular}{p{6.2cm}rrrrrr}
\toprule
Scenario & $R$ & U & Bonferroni & Spec.\ curve & NWC package & NWC H1 \\
\midrule
Null (no effect; treatment exchangeable within sector) & 200 & 0.0 & 0.0 & 11.0 & 0.0 & 5.0 \\
Null, moderators chosen by peeking at outcomes (not replayed) & 300 & 0.0 & 0.0 & 10.0 & 0.0 & 10.3 \\
Null, the same peeking step replayed in every null world & 200 & 0.0 & 0.0 & 11.0 & 0.0 & 5.0 \\
No causal effect; unobserved trait drives both exposure and post-period drift & 50 & 0.0 & 0.0 & 94.0 & 4.0 & 90.0 \\
Planted effects (H1 0.10 SD; three moderators $\pm$0.06 SD) & 50 & 82.0 & 60.0 & 100.0 & 82.0 & 100.0 \\
\bottomrule
\end{tabular}
\begin{minipage}{0.95\linewidth}\vspace{2pt}\scriptsize\textit{Notes.} Each replication draws a synthetic twin (Section~\ref{sec:mc}), runs the reduced registered loop (12 base specifications and the 10-moderator search), and computes loop-level $p$-values from $B$ inner null worlds (within-sector permutation of exposure). Under the null, NWC rejection rates should not exceed 5\% up to Monte Carlo error; Theorem 1 does not cover the leakage and confounded scenarios, which show what the method cannot do.\end{minipage}
\end{table}